%===============================================================
% chap3.tex — Firmware và thuật toán
%===============================================================
\pagestyle{fancy}
\chapter{FIRMWARE VÀ THUẬT TOÁN OSCILLOMETRIC}

\section{Kiến trúc phần mềm nhúng}
Firmware được xây dựng bằng PlatformIO, board \texttt{bluepill\_f103c8}, framework STM32Cube HAL. Các module chính:
\begin{itemize}\itemsep0pt
	\item Driver ADS1115 đọc kênh vi sai phù hợp với điện áp ra của MPX5050 và PGA đặt trong phạm vi đầy thang không bão hòa.
	\item Chuyển đổi đếm ADC \(\rightarrow\) áp suất cuff (mmHg) bằng offset và hệ số scale (hiệu chỉnh thực nghiệm).
	\item FSM đo huyết áp: pha bơm chậm có ``nghe'' dao động; pha ramp đến mục tiêu sau lệnh host; pha xả chậm để thu sóng oscillometric; pha xả nhanh kết thúc hoặc an toàn.
	\item UART text-line: stream mẫu và thông báo trạng thái cho WebApp.
\end{itemize}

\section{Luồng đo và giao thức UART}
Bảng~\ref{tab:uart-frames} tóm tắt các khung điển hình (chi tiết trong tài liệu \texttt{pc-bridge.md} của kho mã).

\begin{table}[htbp]
	\centering
	\caption{Các dòng lệnh/thông báo UART tiêu biểu.}
	\label{tab:uart-frames}
	\begin{tabular}{ll}
		\hline
		\textbf{Dòng} & \textbf{Ý nghĩa} \\
		\hline
		\texttt{S,<seq>,<t\_ms>,<p\_mmHg>} & Stream áp suất (~100 Hz trong đo) \\
		\texttt{A,IDLE} / \texttt{A,INFLATE\_SLOW} / \texttt{A,DEFLATE} & Trạng thái FSM \\
		\texttt{T,<target\_mmHg>} & Host đặt áp mục tiêu (clamp \(\leq\) SAF) \\
		\texttt{ABORT} hoặc \texttt{A} & Yêu cầu xả nhanh an toàn \\
		\texttt{E,MEAS\_END} & Kết thúc đo \\
		\hline
	\end{tabular}
\end{table}

Tốc độ nền: baud 115200, kết thúc dòng \texttt{\textbackslash r\textbackslash n}. Trần SAF \textbf{280~mmHg} áp dụng đồng nhất cho điều khiển cục bộ và clamp lệnh \texttt{T,...}.

\section{Cơ sở thuật toán MAA trên host}
Thuật toán biên độ cực đại (MAA) làm việc trên đường bao dao động (OWE) của thành phần AC của áp suất cuff trong quá trình xả chậm:
\begin{enumerate}\itemsep0pt
	\item Xác định \(\mathrm{MAP}\) là áp suất cuff tại điểm biên độ dao động cực đại \(A_{\max}\).
	\item Với các hệ số đặc trưng \(r_s\) (tâm thu) và \(r_d\) (tâm trương), tính các mức biên độ mục tiêu \(A_{\mathrm{sys}} = r_s A_{\max}\), \(A_{\mathrm{dia}} = r_d A_{\max}\).
	\item Dò nhánh áp cao hơn MAP để tìm điểm gần \(A_{\mathrm{sys}}\) \(\Rightarrow\) SBP; dò nhánh áp thấp hơn MAP để tìm điểm gần \(A_{\mathrm{dia}}\) \(\Rightarrow\) DBP.
\end{enumerate}

Trên host, trước khi trích envelope và cực đại, tài liệu kỹ thuật dự án khuyến nghị \textbf{lọc số dải thông 0{,}5--5~Hz} để giữ nhịp mạch và giảm nhiễu tay/run.

\section{Đồng bộ với pha bơm chậm}
Trong pha \texttt{INFLATE\_SLOW}, WebApp theo dõi dao động sau lọc để ước lượng \(\mathrm{P\_sys\_est}\). Sau đó host gửi \(\texttt{T},\mathrm{P\_sys\_est}+40\) (margin mặc định trong \texttt{board\_config.h}) để MCU ramp đến mục tiêu rồi chuyển xả chậm đo. Nếu không có host trong thời gian chờ, firmware có thể ramp theo mục tiêu dự phòng (\texttt{INFLATE\_FALLBACK\_TARGET\_MMHG}).

\section{LED và trải nghiệm người dùng}
LED phản ánh trạng thái không chặn luồng đo (tham chiếu \texttt{LED\_Algorithm.md}): idle/done xanh; bơm margin vàng nhấp nháy; đo/xả vàng sáng; STOP/SAF/\texttt{ABORT} đỏ nhấp nháy nhanh; lỗi rò/I2C đỏ sáng.
